% Source: https://artofproblemsolving.com/community/c3h3166099_2019_lmt_spring_guts_round_p25p36__lexington_mathematical_tournament
% Stable identity: aops-topic-3166099-q29
The expression equals $1+\cos(x-y)+\cos(y-z)+\cos(z-x)=4\cos\dfrac{x-y}{2}\cos\dfrac{y-z}{2}\cos\dfrac{z-x}{2}$. Order the three angles around the circle. Examining the signs of these half-angle cosines shows the product is positive exactly when one of the three circular gaps exceeds $\pi$, equivalently when the three points lie in an open semicircle. For each chosen point, the probability that both others lie within the following clockwise semicircle is $\dfrac{1}{4}$. These three events are disjoint except for zero-probability boundary cases, since a gap exceeding $\pi$ is unique. Their total probability is $3\cdot\dfrac{1}{4}=\boxed{\dfrac{3}{4}}$.
\begin{asy}
usepackage("amsmath");
size(7cm);
draw(circle((0,0),1));pair[]v={dir(20),dir(70),dir(150)};for(int j=0;j<3;++j){draw((0,0)--v[j],dashed);dot(v[j]);label("$"+string(j+1)+"$",v[j],v[j]);}draw(arc((0,0),1.1,20,200),blue);label("one containing semicircle",(0,-1.4));
\end{asy}
